\section{Clusters: conservative dynamics and observed targets}
\label{sec:clusters}

The cluster calculation uses positive, inertial stock, gas, and stars in
one gravitational potential. Gas pressure and shocks can separate the
gas from the weakly coupled stock and galaxies during an encounter.
The observed lensing map measures their total projected mass. This
construction supplies a transport mechanism and quantitative targets;
the present comparisons do not constitute a fitted cluster population.
\source{CL1--CL3}

\subsection{A common mass budget and a conservative encounter}
For a fully released closed cohort, the adopted normalization is
\begin{equation}
k_C=e^{-\eta_*}/\epsilon=5.1760752,\qquad C_0=k_C B_0.
\end{equation}
If its baryons and stock remain coextensive, the baryon fraction is
$1/(1+k_C)=0.1619151$. More generally, an aperture retains different
fractions of the ancestral components. The measured target can be
expressed as
\begin{equation}
C_{\rm req}=M_{\rm grav}-B,\qquad
\eta_{\rm ret}=\frac{k_C B}{C_{\rm req}}
=\frac{f_b^{\rm ret}}{f_C^{\rm ret}}.
\end{equation}
This distinguishes the common cohort normalization from the distribution
produced by assembly. A positive inventory exists for $C_{\rm req}>0$:
$B_0=\max(B,C_{\rm req}/k_C)$ leaves nonnegative exterior baryons
$B_0-B$ and stock $k_CB_0-C_{\rm req}$. It does not determine a
realizable transport history or its energy cost.

The encounter model uses central pair impulses with opposite momentum
changes. If $w$ is the radial relative velocity and $\mu$ the reduced
mass, the kinetic change is $\Delta K=\mu(w'^2-w^2)/2$. Explicit work
and heat registers balance that change, and a symmetric gate suppresses
transitions with an endpoint at the energy floor. Gas and stars from
either progenitor can be local partners. The target distribution, rate,
range, and initial work reserves are selected effective inputs.
\source{D2, CL1}

Fourteen encounters and three isolated controls test a common Plummer
mass ratio of 0.3, with 1344 particles at baseline and 2688 in the
resolution control. Weak contact leaves stock near the galaxies while
gas lags. A floor-state run permits no events and reproduces the
no-contact evolution; a low-work run permits 61 events against 7114
with the larger reserve. The largest energy drift is below 0.007\%
of the binding energy; pair momentum conservation holds at roundoff.
The strongest tested contact changes the isolated stock median radius
by $-27\%$ and gas median radius by $+46\%$ in six dynamical time units.
That control prevents interpreting a higher contact rate as an established
cluster solution. The particle encounters test the finite-work coupling;
they do not evolve the full spatial stress of the depleted scalar.

\subsection{X-COP profiles and independent weak-lensing masses}
The public X-COP archive supplies twelve hydrostatic radial profiles,
seven with stellar mass data. Five overlap the MENeaCS lensing catalogue.
The comparison adopts the observational distance convention
$H_0=70\,\kms\,\mathrm{Mpc}^{-1}$ and $\Omega_m=0.3$.
Projected stellar masses are approximately deprojected using the stated
profile model. Unmeasured intracluster light, hydrostatic departures,
triaxiality, and profile covariance remain relevant.
\cite{xcopnew,herbonnetnew}\source{CL2}

All nominal NFW hydrostatic targets admit positive cumulative stock and
positive stock in fifteen tested shells. At the hydrostatic $R_{500}$,
the coextensive-to-hydrostatic total-mass ratios for A1795, A2029, A2142,
A2319, A644, A85, and ZW1215 are respectively
$0.903$, $0.912$, $1.023$, $1.213$, $0.851$, $0.952$, and $0.679$.
These reconstructed totals are inputs to a target audit, rather than
stock profiles predicted from the observed baryons.

\begin{table}[htbp]
\centering\small
\begin{tabular}{lrrr}
\toprule
Cluster & Coextensive $M_{500}$ & Lensing $M_{500}$ & $\eta_{\rm ret}$\\
\midrule
A1795 & 3.96--4.23 & $9.3\pm2.2$ & 0.518\\
A2029 & 7.46--8.07 & $12.1\pm2.5$ & 0.704\\
A2142 & 9.29--10.02 & $9.7\pm2.3$ & 0.969\\
A85 & 5.20--5.43 & $5.7\pm2.2$ & 0.937\\
ZW1215 & 4.19--4.46 & $3.5\pm2.2$ & 1.143\\
\bottomrule
\end{tabular}
\caption{Masses in $10^{14}\Msun$. The coextensive prescription solves
its own overdensity radius; its range doubles the nominal stellar mass
and is a sensitivity interval. Retention ratios use the nominal baryons
within each central lensing target's radius. Lensing errors are the
published NFW marginal errors. No combined exclusion significance is
assigned.}
\end{table}

\begin{figure}[htbp]
\centering\includegraphics[width=\textwidth]{figures/cluster_lensing.png}
\caption{The common coextensive prescription and two mass estimators
from the same lensing data. Their errors are correlated. A1795 and
A2029 require a transport or reconstruction explanation; A2142 and A85
are useful normalization controls. ZW1215 is sensitive to the estimator.}
\end{figure}

\subsection{What the outliers require}
For central lensing targets, the minimal A1795 cohort leaves 48.2\%
of its initial baryons outside the aperture, or
$7.87\times10^{13}\Msun$; A2029 leaves 29.6\% outside. Twelve saved
finite transport maps preserve nonnegative marginals for the seven
hydrostatic and five lensing targets. They establish inventory
compatibility. They do not establish that the selected coupling produces
the required segregation in the available time. \source{CL3}

Gas support must use the same total potential. In a radial description,
\begin{equation}
a_g=-g-\frac{P_{\rm th}'+P_{\rm nt}'}{\rho_g}
       +\frac{f_{Cg}}{\rho_g}.
\end{equation}
Assembly segregation, gas acceleration, and reciprocal contact are
possible contributions to the discrepancy. A passive drag vanishes at
zero relative drift, so a stationary contact force requires a separate
physical law. XRISM measurements give low turbulent support in A1795
and A2029, constraining a purely turbulent explanation.
\cite{sarkarnew,xrismnew} A conditional pressure-integral audit of the
central lensing profiles also requires the inner extrapolation and
thermal reconstruction to be propagated; it does not provide a joint
exclusion or a completed contact-force fit.

A644 has a distinct reconstruction issue: its 58-knot forward profile
contains 11 negative total-mass increments and 15 negative inferred-stock
increments. A monotone profile requires a maximum displacement of
2.156 tabulated marginal error units. The unknown joint covariance
precludes treating this as a Gaussian exclusion significance. Positive
parametric mass profiles are available and must be tested against the
underlying pressure and density data. For ZW1215, at 1320 kpc the NFW
lensing, aperture-lensing, and hydrostatic masses are respectively
$4.32$, $7.0$, and $7.41\times10^{14}\Msun$. That estimator spread must
be retained in a physical interpretation. These systems remain open
tests of a common dynamics, with no cluster-specific gravity coefficient.

\subsection{The Bullet Cluster supplies a spatial target}
The CANUCS Bullet v1 best-fit maps separate a total lens reconstruction
from a supplied gas component. At lens redshift 0.296 and
$D_{ls}/D_s=1$, their normalization gives
$\Sigma_{\rm crit}=1.8268761\times10^9\Msun\,\mathrm{kpc}^{-2}$
and $1''=4.41302$ kpc. Registered apertures around BCG3 give
\begin{center}
\begin{tabular}{lrrr}
\toprule
Radius & $M_{\rm total}/10^{13}\Msun$ & $M_g/10^{13}\Msun$ & Gas fraction\\
\midrule
60 kpc & 2.151 & 0.099 & 0.046\\
150 kpc & 8.511 & 0.618 & 0.073\\
\bottomrule
\end{tabular}
\end{center}
The released subcluster halo centre lies 3.92 kpc from BCG3. The
supplied gas-density map peaks about 169 kpc away within the specified
search region, with an 8.68 kpc gas pixel. This is a density-map
diagnostic, rather than an X-ray brightness peak or an offset posterior.
Subtracting gas leaves a component that still includes stars.
\cite{bulletnew}\source{CL2}

\begin{figure}[htbp]
\centering\includegraphics[width=\textwidth]{figures/cluster_bullet.png}
\caption{Observed targets from the released Bullet v1 lens model.
The red contours show its supplied gas component. Both panels are
derived from the public reconstruction; no simulated map is fitted here.}
\end{figure}

Aperture masses change by at most 0.214\% under pixel coarsening; this
is a numerical check and excludes observational uncertainty. Under an
illustrative $10^{15}\Msun$, 300 kpc scaling, the encounter calculations
have 24--30 kpc force softening and roughly 90 kpc map smoothing. They
cannot establish agreement with the few-kiloparsec halo--galaxy alignment.
A successful forward test must preserve isolated cluster profiles and
then reproduce lensing, gas morphology, and offsets using one selected
contact law and independently specified assembly conditions.

